% READCRAFTERY — Artist Bible v1.2
% Compile with: xelatex READCRAFTERY_ARTIST_BIBLE_v1_2.tex
% Images must be in the same directory as this .tex file

\documentclass[11pt, a4paper]{article}

\usepackage{fontspec}
\usepackage{geometry}
\usepackage{xcolor}
\usepackage{graphicx}
\usepackage{booktabs}
\usepackage{tabularx}
\usepackage{array}
\usepackage{colortbl}
\usepackage{enumitem}
\usepackage{titlesec}
\usepackage{fancyhdr}
\usepackage{hyperref}
\usepackage{tikz}
\usepackage{mdframed}
\usepackage{parskip}
\usepackage{microtype}
\usepackage{float}

\geometry{a4paper, top=2cm, bottom=2cm, left=2.2cm, right=2.2cm, headheight=14pt}

% Colors
\definecolor{navycol}{HTML}{1E2A4A}
\definecolor{goldcol}{HTML}{D4A017}
\definecolor{parchmentcol}{HTML}{F5E6C8}
\definecolor{parchment2col}{HTML}{E8D4A8}
\definecolor{browncol}{HTML}{2C1A0E}
\definecolor{surfacecol}{HTML}{4A2E1A}
\definecolor{magiccol}{HTML}{7B4FBE}
\definecolor{correctcol}{HTML}{4CAF50}
\definecolor{incorrectcol}{HTML}{E53935}
\definecolor{neutralcol}{HTML}{9E9E9E}
\definecolor{lightnote}{HTML}{FFF8E7}
\definecolor{warnbg}{HTML}{FFF3E0}
\definecolor{warnring}{HTML}{E65100}
\definecolor{tabhead}{HTML}{1E2A4A}
\definecolor{tabrow}{HTML}{F9F6F0}
\definecolor{m2color}{HTML}{C62828}
\definecolor{m3color}{HTML}{E65100}
\definecolor{m2bg}{HTML}{FFEBEE}
\definecolor{m3bg}{HTML}{FFF3E0}
\definecolor{okbg}{HTML}{E8F5E9}
\definecolor{okcolor}{HTML}{1B5E20}
\definecolor{storytext}{HTML}{3A2510}

\setmainfont{DejaVu Serif}
\setsansfont{DejaVu Sans}
\setmonofont{DejaVu Sans Mono}

\hypersetup{colorlinks=true, linkcolor=navycol, urlcolor=magiccol,
  pdftitle={READCRAFTERY Artist Bible v1.2}, pdfauthor={The Readcraftery Project}}

\titleformat{\section}{\sffamily\Large\bfseries\color{navycol}}{\thesection.}{0.5em}{}
  [\color{goldcol}\titlerule[1.5pt]\vspace{2pt}]
\titleformat{\subsection}{\sffamily\large\bfseries\color{navycol}}{\thesubsection}{0.5em}{}
\titleformat{\subsubsection}{\sffamily\normalsize\bfseries\color{surfacecol}}{}{0em}{}
\titlespacing*{\section}{0pt}{18pt}{8pt}
\titlespacing*{\subsection}{0pt}{12pt}{4pt}
\titlespacing*{\subsubsection}{0pt}{8pt}{2pt}

\pagestyle{fancy}
\fancyhf{}
\fancyhead[L]{\small\sffamily\color{navycol}\textbf{READCRAFTERY} -- Artist Bible v1.2}
\fancyhead[R]{\small\sffamily\color{neutralcol}March 2026}
\fancyfoot[C]{\small\sffamily\color{neutralcol}\thepage}
\renewcommand{\headrulewidth}{0.4pt}

% Note box
\newmdenv[backgroundcolor=lightnote, linecolor=goldcol, linewidth=2pt,
  leftline=true, rightline=false, topline=false, bottomline=false,
  innerleftmargin=10pt, innerrightmargin=8pt, innertopmargin=6pt, innerbottommargin=6pt,
  skipabove=8pt, skipbelow=8pt]{notebox}

% Warning box
\newmdenv[backgroundcolor=warnbg, linecolor=warnring, linewidth=2pt,
  leftline=true, rightline=false, topline=false, bottomline=false,
  innerleftmargin=10pt, innerrightmargin=8pt, innertopmargin=6pt, innerbottommargin=6pt,
  skipabove=8pt, skipbelow=8pt]{warnbox}

% Story box
\newmdenv[backgroundcolor=parchmentcol!25, linecolor=browncol!35, linewidth=1pt,
  leftline=true, rightline=true, topline=true, bottomline=true,
  innerleftmargin=14pt, innerrightmargin=14pt, innertopmargin=10pt, innerbottommargin=10pt,
  skipabove=10pt, skipbelow=10pt, roundcorner=4pt]{storybox}

% \hex{RRGGBB} — shows hex code in its own color
\newcommand{\hex}[1]{\textcolor[HTML]{#1}{\texttt{\##1}}}

% \colorswatch{colorname}{RRGGBB}{description}
\newcommand{\colorswatch}[3]{%
  \noindent
  \begin{tikzpicture}[baseline=-0.5ex]
    \fill[#1] (0,0) rectangle (0.7cm, 0.35cm);
    \draw[gray!40] (0,0) rectangle (0.7cm, 0.35cm);
  \end{tikzpicture}
  \hspace{4pt}\hex{#2}\hspace{8pt}{\sffamily\small #3}\\[3pt]
}

% Priority badge
\newcommand{\badgeM}[1]{%
  \colorbox{\ifnum#1=2 m2bg\else\ifnum#1=3 m3bg\else okbg\fi\fi}{%
    \color{\ifnum#1=2 m2color\else\ifnum#1=3 m3color\else okcolor\fi\fi}%
    \sffamily\small\bfseries M#1}%
}

% ─────────────────────────────────────────────────────────────────────────────
\begin{document}
\thispagestyle{empty}

% COVER
\begin{center}
  \vspace*{1.5cm}
  {\fontsize{42}{50}\selectfont\sffamily\bfseries\color{navycol} READCRAFTERY}\\[6pt]
  {\Large\sffamily\color{surfacecol} Artist Bible --- Visual Canon and Production Rules}\\[4pt]
  {\normalsize\sffamily\color{neutralcol}\itshape Version 1.1 --- March 2026}\\[2cm]
  \includegraphics[width=7cm]{owlorumo_concept}\\[6pt]
  {\small\sffamily\itshape\color{neutralcol} Owlorumo --- Official Concept Art}\\[1.5cm]
  \begin{mdframed}[linecolor=goldcol, linewidth=1.5pt,
      innerleftmargin=16pt, innerrightmargin=16pt,
      innertopmargin=8pt, innerbottommargin=8pt]
    {\small\sffamily\itshape\color{neutralcol} Confidential --- Not for distribution}
  \end{mdframed}
\end{center}

\newpage
\setcounter{page}{1}
{\sffamily\tableofcontents}
\newpage

% ─────────────────────────────────────────────────────────────────────────────
\section{Project Overview}

READCRAFTERY is an educational word puzzle game for children ages 5--9, built in Godot 4.
The game teaches reading through interactive puzzles set in a magical library world.
Players read short passages, then search for target words guided by their companion
Owlorumo --- a young owl wizard whose power returns as the child reads.

The visual style is pixel art with a warm, magical, child-friendly aesthetic:
\textbf{Stardew Valley} meets a \textbf{1980s illustrated storybook}.

\subsection{Target Platforms}
\begin{itemize}[leftmargin=1.5em, itemsep=2pt]
  \item HTML5 (itch.io) --- primary
  \item Android / iOS
  \item PC (Windows / Linux / macOS)
\end{itemize}

\subsection{Art Style --- Storybook Arcane}

The style has a name: \textbf{Storybook Arcane}. Key principles:

\begin{itemize}[leftmargin=1.5em, itemsep=2pt]
  \item Pixel art with deliberate organic shapes --- curves, not sharp geometry
  \item Warm candlelight as primary light source --- upper left in all scenes
  \item Magic is ambient, not explosive --- soft glows, floating particles, pulsing runes
  \item Textures are invented --- textures that only exist in this world
  \item Characters have exaggerated expressiveness --- big eyes, never frightening
  \item No dark fantasy --- no skulls, no scary darkness; the wizard is always kind
\end{itemize}

\subsection{The Library Through Restoration}

The same place is encountered under three progressively revealed names:
\begin{itemize}[leftmargin=1.5em, itemsep=2pt]
  \item \textbf{The Sleeping Library} --- the first state; dormant, veiled, waiting
  \item \textbf{The Reawakened Library} --- a warmer, more responsive state after meaningful restoration
  \item \textbf{The Library of Beginnings} --- the deeper revealed truth; restoration does not return a pristine past, it opens new continuity
\end{itemize}

These are not separate worlds or separate levels. They are names for the same living library-world as it becomes more present.

\subsection{Master Color Palette}

All assets must use \textbf{only} these colors. No additional colors without approval.

\vspace{6pt}
\colorswatch{browncol}{2C1A0E}{background --- main background, dark warm brown}
\colorswatch{surfacecol}{4A2E1A}{surface --- shelves, furniture, elevated surfaces}
\colorswatch{parchmentcol}{F5E6C8}{parchment --- reading area background, light UI}
\colorswatch{parchment2col}{E8D4A8}{parchment\_shadow --- secondary text backgrounds}
\colorswatch{goldcol}{D4A017}{accent (gold) --- restoration cues, warm highlights, identity elements}
\colorswatch{magiccol}{7B4FBE}{magic\_glow --- staff crystal, rune glows, puzzle highlights}
\colorswatch{correctcol}{4CAF50}{correct --- correct answer feedback}
\colorswatch{incorrectcol}{E53935}{incorrect --- use sparingly, never harsh}
\colorswatch{neutralcol}{9E9E9E}{neutral --- disabled states, dormant content}

\begin{warnbox}
  \textbf{Warning:} \texttt{magic\_glow} (\hex{7B4FBE}) appears \textit{only} on Owlorumo's
  staff crystal, puzzle highlight effects, and narrative rune glows. Never in functional UI.\\[4pt]
  The accent gold (\hex{D4A017}) is sacred --- do not substitute with other yellows.
  Gold should feel meaningful and warm, but never like score or trophy UI.
\end{warnbox}

% ─────────────────────────────────────────────────────────────────────────────
\section{Owlorumo --- Character Design}

\subsection{The Story (context for the artist)}

\begin{storybox}
{\sffamily\small\color{storytext}
Owlorumo was the companion owl of Athena. As the habit of reading dimmed in the human world,
his memory, power, and presence diminished with it. He remains in the library --- the last
place where wisdom survived long enough to be reawakened.

\vspace{6pt}
The child is never told this directly. Owlorumo does not explain it. But each recognized word,
pattern, and finished passage returns a little more light to what was fading.

\vspace{6pt}
\textbf{Reading is a power that discovers and restores human knowledge.}

\vspace{6pt}
\textit{This story lives in visual details: the staff crystal, Owlorumo's eyes, the warmth of
his presence, and the awakening of the library itself. \textbf{state\_4} is a deliberate homage
to Tolkien's Gandalf the White: fulfilled, not replaced. The blue mantle remains visible.
The moon brooch, belt, boots, and organic staff form remain. The transformation is one of
light, material, and presence. Never explained. Only discovered.}
}
\end{storybox}

\subsection{Personality}
\begin{itemize}[leftmargin=1.5em, itemsep=2pt]
  \item Young, curious, slightly clumsy with magic but never with words
  \item Companion, not authority --- discovers alongside the child
  \item Amber eyes are his most expressive feature: wide with wonder, warm with encouragement
  \item Never intimidating, never scared, never disappointed
  \item No beard --- Owlorumo is young
\end{itemize}

\subsection{Reference Images}

The following two images are the \textbf{official design reference}.
All pixel art must match these. When in doubt, the concept art wins.

\vspace{8pt}
\begin{figure}[H]
  \centering
  \begin{minipage}{0.46\textwidth}
    \centering
    
\includegraphics[width=\linewidth]{../assets/img/owlorumo_concept_art}\\
    {\small\itshape\color{neutralcol} Concept Art --- official reference}
  \end{minipage}
  \hfill
  \begin{minipage}{0.46\textwidth}
    \centering
    \includegraphics[width=\linewidth]{owlorumo_concept_pixel_art}\\
    {\small\itshape\color{neutralcol} AI Pixel Art --- style guide}
  \end{minipage}
\end{figure}

\subsection{Character Anatomy}

\subsubsection{Proportions}
Compact and slightly chubby --- reads as young. Large round head. \textbf{No beard.}

\subsubsection{Clothing --- confirmed from concept art}
\begin{itemize}[leftmargin=1.5em, itemsep=2pt]
  \item Navy blue robe --- deep, dark (\hex{1E2A4A}). Slightly oversized (apprentice feel)
  \item Gold trim on robe hem and cuffs (\hex{D4A017})
  \item Neck wrap --- same navy fabric, loosely wrapped
  \item \textbf{Crescent moon brooch} --- \hex{D4A017}. Always present. Identity detail.
  \item Brown leather belt with gold buckle
  \item Tall pointed wizard hat --- same navy, slightly bent. Feather + rune engravings.
  \item Short brown leather boots --- slightly oversized, covers talons completely
\end{itemize}

\subsubsection{Staff --- Integral Element (NEVER remove)}

The staff is a \textbf{permanent} part of Owlorumo's design --- present in every frame.

\begin{itemize}[leftmargin=1.5em, itemsep=2pt]
  \item Dark twisted organic wood --- darker than the robe
  \item Purple crystal on top --- \hex{7B4FBE} in idle, brighter in active states
  \item Owlorumo holds the staff in his \textbf{LEFT hand}, reaching above his head
  \item The crystal glow is the single most recognizable silhouette element after the hat
\end{itemize}

\begin{notebox}
  Owlorumo's silhouette must always read as ``owl with wizard hat and staff'' at 64x64px.
  If the staff cannot fit legibly, show the crystal glow at minimum.
\end{notebox}

\subsection{Eyes}

At 64x64px the eyes will be 4--5 pixels wide each. Every emotion comes through the eyes.

\begin{itemize}[leftmargin=1.5em, itemsep=2pt]
  \item Bright amber-gold (\hex{D4A017} tint toward orange)
  \item Large and round --- Seuss-style, not realistic
  \item Independent eyebrows (owls don't have them --- intentional)
  \item Slow blink in idle = signal of trust and warmth
  \item \textbf{Each eye has a 1-pixel white highlight in the upper-left corner.}
        This is what makes them read as alive. Do not omit.
\end{itemize}

\subsection{Companion Progression States}

Owlorumo starts diminished and is progressively restored as the child reads.
\textbf{The artist delivers 5 PNG atlas files with identical frame layout.}

\vspace{6pt}
\begin{tabularx}{\textwidth}{@{}c >{\ttfamily\small}l >{\raggedright\arraybackslash}X
    >{\raggedright\arraybackslash}X >{\raggedright\arraybackslash}X@{}}
  \toprule
  \rowcolor{tabhead}
  \color{white}\sffamily\bfseries \# &
  \color{white}\sffamily\bfseries File &
  \color{white}\sffamily\bfseries Crystal &
  \color{white}\sffamily\bfseries Robe &
  \color{white}\sffamily\bfseries Unlocked after \\
  \midrule
  \rowcolor{tabrow}0 & atlas\_owlorumo\_state0.png & Cracked, dark & Navy, no runes & Game start \\
  1 & atlas\_owlorumo\_state1.png & Faint \hex{7B4FBE} glow & Navy, no runes & First passage \\
  \rowcolor{tabrow}2 & atlas\_owlorumo\_state2.png & Lit, pulsing & Navy, first rune & First book \\
  3 & atlas\_owlorumo\_state3.png & Bright \hex{7B4FBE} & Navy, several runes & 3 books \\
  \rowcolor{tabrow}4 & atlas\_owlorumo\_state4.png & \textbf{White-gold, blinding} & \textbf{Blue mantle illuminated white-gold; same brooch, belt, boots} & First complete arc \\
  \bottomrule
\end{tabularx}

\vspace{6pt}
\begin{notebox}
  \textbf{Crystal behavioral states — animation guidance:}\\[4pt]
  The crystal communicates four distinct states through animation:
  \begin{itemize}[leftmargin=1.5em, itemsep=2pt]
    \item \textbf{Idle pulse} — soft \hex{7B4FBE} pulse at 0.8Hz. Base state.
    \item \textbf{Hint active} — brightens to \hex{A87FE8}, staff points toward puzzle area.
    \item \textbf{Fragment absorption} — brief white-gold flash when a puzzle is completed.
      The flash arrives, pulses once, is absorbed. Duration: under 1 second.
      Visually distinct from the idle pulse — brighter, warmer, directional.
    \item \textbf{State base level} — permanent ambient glow determined by owlorumo\_state.
      This is the resting brightness between all other states.
      Fragments arrive on top of this level — they do not replace it.
  \end{itemize}
  The artist must design the crystal so all four states are readable
  at 64×64px native size.
\end{notebox}
\begin{notebox}
  \textbf{state\_0 and state\_1 are M2 blockers} --- deliver these first.\\[4pt]
  \textbf{state\_4} is a homage to Tolkien's Gandalf the White. Never explained in the game.
  Only discovered. The crystal in state\_4 transitions from \hex{7B4FBE} to a white-gold glow
  --- the only state where the crystal color changes.
\end{notebox}

\subsection{Required Animation States}

Deliver as sprite atlas PNG, transparent background, \textbf{64x64px per frame}.
All 5 state atlases share the same frame layout.

\vspace{6pt}
\begin{tabularx}{\textwidth}{@{}l c l >{\raggedright\arraybackslash}X@{}}
  \toprule
  \rowcolor{tabhead}
  \color{white}\sffamily\bfseries State &
  \color{white}\sffamily\bfseries Frames &
  \color{white}\sffamily\bfseries Priority &
  \color{white}\sffamily\bfseries Description \\
  \midrule
  \rowcolor{tabrow}\texttt{owl\_idle}      & 4 fr & \badgeM{2} & Soft breathing. Slow blink fr 3. Staff crystal pulses 0.8Hz. \\
  \texttt{owl\_happy}     & 4 fr & \badgeM{2} & Neck feathers ruffled. Slight sway. Half-closed eyes. \\
  \rowcolor{tabrow}\texttt{owl\_hint}      & 4 fr & \badgeM{2} & Leans forward. Staff gently points toward puzzle. Crystal bright. \textbf{Owlorumo is the hint button — the child taps him directly.} \\
  \texttt{owl\_thinking}  & 4 fr & \badgeM{2} & Slow head turn. Talon scratches chin. Staff in other hand. \\
  \rowcolor{tabrow}\texttt{owl\_excited}   & 6 fr & \badgeM{3} & Wings slightly extended. Short hop. Crystal flashes. \\
  \texttt{owl\_celebrate} & 8 fr & \badgeM{3} & Wings fully extended. Magic motes emerge from crystal. \\
  \rowcolor{tabrow}\texttt{owl\_reading}   & 4 fr & \badgeM{3} & Holds open book. Eyes move side to side slowly. \\
  \texttt{owl\_sleepy}    & 4 fr & \badgeM{3} & Slow nodding. Half-closed eyes. Inactivity screens. \\
  \rowcolor{tabrow}\texttt{owl\_wave}      & 4 fr & \badgeM{3} & Raises one wing in greeting. Onboarding and return. \\
  \bottomrule
\end{tabularx}

\vspace{6pt}
\begin{warnbox}
  \textbf{M2 = first 4 states are blockers.}
  Deliver \texttt{owl\_idle}, \texttt{owl\_happy}, \texttt{owl\_hint},
  \texttt{owl\_thinking} for \textbf{state\_0} and \textbf{state\_1} first.\\[4pt]
  Display is controlled by Godot per context:
  192x192px in PassageView (3x), 256x256px in Library/menu (4x).
  Both are integer scales --- no quality loss.
\end{warnbox}

\subsection{Owlorumo in the Library — Ambient Behavior}

Owlorumo is not static in the library. He inhabits the space.
The artist must design animations that communicate this.

\begin{notebox}
  \textbf{Core rule:} Owlorumo is always doing something when the child arrives.
  He finishes his current gesture naturally, then looks at the child.
  He was not waiting. He was here. There is a difference.
\end{notebox}

\textbf{Ambient activities — same action, different energy per state:}

\begin{tabularx}{\textwidth}{@{}l >{\raggedright\arraybackslash}X >{\raggedright\arraybackslash}X@{}}
  \toprule
  \rowcolor{tabhead}
  \color{white}\sffamily\bfseries Activity &
  \color{white}\sffamily\bfseries state\_0 / state\_1 &
  \color{white}\sffamily\bfseries state\_3 / state\_4 \\
  \midrule
  \rowcolor{tabrow}Sorting books with magic & Slow, a book falls, he picks it up & Fluid, effortless \\
  Walk to window & Heavy, slightly hunched & Light, curious \\
  \rowcolor{tabrow}Reading at lectern & Eyes moving slowly, sometimes stops & Engaged, occasional soft smile \\
  Writing on parchment & Only state\_2+ & Present from state\_2 \\
  \bottomrule
\end{tabularx}

\vspace{6pt}
\textbf{On child entry:} Owlorumo always notices. He finishes his gesture,
then turns toward the child. Never startled. Never ignoring.

\textbf{On book selection:} a small gesture of recognition —
not celebration, not instruction. As if the child chose exactly right.

\textbf{On tap in library (no active event):} crystal pulses softly,
a glance toward the child. Presence, not function.

% ─────────────────────────────────────────────────────────────────────────────
\section{Environment Assets}

All backgrounds: \textbf{draw at 480x270px}, displayed at 1920x1080 (4x).
Light source always upper left.

\subsection{Library Scene}

\begin{figure}[H]
  \centering
  \includegraphics[width=0.85\textwidth]{../assets/img/library_scene_concept_art}\\
  {\small\itshape\color{neutralcol} Library Scene --- AI reference (official style reference)}
\end{figure}

Key elements to preserve:
\begin{itemize}[leftmargin=1.5em, itemsep=2pt]
  \item Isometric corner view --- two walls of shelves meeting at center
  \item Candelabra as primary warm light source (upper left)
  \item Purple glowing books on shelves --- \hex{7B4FBE}, ambient not explosive
  \item Reading podium in center. Small desk with quill, ink pot, candle.
  \item Warm stone floor with amber candlelight reflection --- \textbf{not cold grey}
  \item Owlorumo's lectern or nearby presence anchor --- implies he inhabits this space
\end{itemize}

\vspace{6pt}
\begin{tabularx}{\textwidth}{@{}l c c c >{\raggedright\arraybackslash}X@{}}
  \toprule
  \rowcolor{tabhead}
  \color{white}\sffamily\bfseries Asset &
  \color{white}\sffamily\bfseries Draw &
  \color{white}\sffamily\bfseries Display &
  \color{white}\sffamily\bfseries Priority &
  \color{white}\sffamily\bfseries Content \\
  \midrule
  \rowcolor{tabrow}\texttt{bg\_library\_layer\_4} & 480x270 & 1920x1080 & \badgeM{2} & Far background --- dark shelves, stone wall \\
  \texttt{bg\_library\_layer\_3} & 480x270 & 1920x1080 & \badgeM{2} & Mid shelves --- books, some glowing \hex{7B4FBE} \\
  \rowcolor{tabrow}\texttt{bg\_library\_layer\_2} & 480x270 & 1920x1080 & \badgeM{3} & Foreground --- podium, desk, candelabra \\
  \texttt{bg\_library\_layer\_1} & 480x270 & 1920x1080 & \badgeM{3} & Closest --- particles, dust motes \\
  \bottomrule
\end{tabularx}

\begin{figure}[H]
  \centering
  \includegraphics[width=0.65\textwidth]{../assets/img/bookshelf_close_up}\\
  {\small\itshape\color{neutralcol} Bookshelf close-up --- use for layer\_3 detail level}
\end{figure}

\subsection{Library — Diegetic UI elements}

The library UI is part of the world, not drawn on top of it.

\begin{tabularx}{\textwidth}{@{}l >{\raggedright\arraybackslash}X >{\raggedright\arraybackslash}X@{}}
  \toprule
  \rowcolor{tabhead}
  \color{white}\sffamily\bfseries Element &
  \color{white}\sffamily\bfseries Visual identity &
  \color{white}\sffamily\bfseries Position \\
  \midrule
  \rowcolor{tabrow}\texttt{book\_settings} &
    Closed book. Navy \hex{1E2A4A} spine. Crescent moon symbol.
    Visually distinct from all other books. &
    Fixed top-right. Always visible. \\
  \texttt{book\_tts} &
    Small open book. Parchment \hex{F5E6C8}.
    Pulses softly while audio plays. &
    Passage panel, fixed. \\
  \rowcolor{tabrow}\texttt{btn\_exit} &
    Small door. Warm wood. Visible glow on hold.
    Tap-hold 1 second to activate. &
    Bottom-left of PassageView. Fixed. \\
  \bottomrule
\end{tabularx}

\begin{warnbox}
  \textbf{book\_settings color and moon symbol are brand invariants.}
  A theme may restyle the texture but must preserve \hex{1E2A4A} and the crescent moon.
  These are the child's and teacher's navigation anchor — they must be immediately
  recognizable across all themes and sessions.
\end{warnbox}

\textbf{Library books — three visual states:}
\begin{itemize}[leftmargin=1.5em, itemsep=2pt]
  \item \textbf{Available} — saturated colors, subtle gold shimmer, breathes softly.
    Brightens on hover. This is the book that invites.
  \item \textbf{Dormant} — desaturated, no reaction to hover, no lock icon.
    Simply does not respond. Some visible from the first session; others appear
    after completing a passage — placed silently, never announced.
  \item \textbf{Completed} — warm permanent glow. Owlorumo occasionally
    looks toward it during idle.
\end{itemize}

\subsection{PassageView --- Reading Panel}

The reading area. Text legibility is the absolute priority.

\begin{figure}[H]
  \centering
  
\includegraphics[width=0.52\textwidth]{../assets/img/passage.jpg}\\
  {\small\itshape\color{neutralcol} Parchment frame reference --- AI concept}
\end{figure}

\begin{itemize}[leftmargin=1.5em, itemsep=2pt]
  \item Full-panel frame, \textbf{not} a repeating tile
  \item Center must be \textbf{clean} --- no texture, no detail --- text goes here
  \item Corner decorations: crescent moon, quill, book, and celestial glyph motifs --- small, \hex{D4A017}
  \item Faint paper grain. Subtle warm glow from upper left.
\end{itemize}

\begin{tabularx}{\textwidth}{@{}l c c c >{\raggedright\arraybackslash}X@{}}
  \toprule
  \rowcolor{tabhead}
  \color{white}\sffamily\bfseries Asset &
  \color{white}\sffamily\bfseries Draw &
  \color{white}\sffamily\bfseries Display &
  \color{white}\sffamily\bfseries Priority &
  \color{white}\sffamily\bfseries Notes \\
  \midrule
  \rowcolor{tabrow}\texttt{parchment\_tile} & 64x64 & 256x256 & \badgeM{2} & Full panel frame. Clean center. Base \hex{F5E6C8}. \\
  \texttt{wood\_tile} & 64x64 & 256x256 & \badgeM{2} & Puzzle panel. Dark warm wood. Base \hex{4A2E1A}. \\
  \bottomrule
\end{tabularx}

% ─────────────────────────────────────────────────────────────────────────────
\section{UI Assets}

\subsection{Interactive Tiles}

\begin{tabularx}{\textwidth}{@{}l c c c >{\raggedright\arraybackslash}X@{}}
  \toprule
  \rowcolor{tabhead}
  \color{white}\sffamily\bfseries Asset &
  \color{white}\sffamily\bfseries Draw &
  \color{white}\sffamily\bfseries Display &
  \color{white}\sffamily\bfseries Priority &
  \color{white}\sffamily\bfseries Description \\
  \midrule
  \rowcolor{tabrow}\texttt{tile\_word\_normal}   & 16x8 & 64x32 & \badgeM{2} & Word tile. Base \hex{F5E6C8}. Resting state. \\
  \texttt{tile\_word\_selected} & 16x8 & 64x32 & \badgeM{2} & \hex{D4A017} highlight. Pulse animation. \\
  \rowcolor{tabrow}\texttt{tile\_word\_found}    & 16x8 & 64x32 & \badgeM{2} & \hex{4CAF50}. Completed. No animation. \\
  \texttt{slot\_word\_empty}    & 16x8 & 64x32 & \badgeM{2} & Anonymous target slot. \hex{2C1A0E}. Pulse. \\
  \rowcolor{tabrow}\texttt{slot\_word\_filled}   & 16x8 & 64x32 & \badgeM{2} & Slot with found word. \hex{4CAF50}. \\
  \texttt{tile\_letter\_normal} & 16x16 & 64x64 & \badgeM{2} & Square letter tile for Scramble puzzle. \\
  \bottomrule
\end{tabularx}

\subsection{Icons \& Buttons}

\begin{tabularx}{\textwidth}{@{}l c c c >{\raggedright\arraybackslash}X@{}}
  \toprule
  \rowcolor{tabhead}
  \color{white}\sffamily\bfseries Asset &
  \color{white}\sffamily\bfseries Draw &
  \color{white}\sffamily\bfseries Display &
  \color{white}\sffamily\bfseries Priority &
  \color{white}\sffamily\bfseries Description \\
  \midrule
  \rowcolor{tabrow}\texttt{fragment\_light\_01} & 8x8   & 32x32 & \badgeM{2} & Small white-gold light fragment. Idle/shimmer state. \\
  \texttt{fragment\_light\_02} & 8x8   & 32x32 & \badgeM{2} & Alternate fragment shape for travel/arrival variation. \\
  \rowcolor{tabrow}\texttt{icon\_hint}      & 16x16 & 64x64 & \badgeM{2} & Hint button --- staff crystal motif. \hex{7B4FBE}. \\
  \texttt{btn\_back}       & 24x24 & 96x96 & \badgeM{2} & Back/home button. Normal + pressed states. \\
  \rowcolor{tabrow}\texttt{book\_available} & 16x24 & 64x96 & \badgeM{2} & Book cover for Library Scene. \\
  \texttt{book\_dormant}   & 16x24 & 64x96 & \badgeM{2} & Dormant book. Desaturated, quiet, no lock icon. \\
  \bottomrule
\end{tabularx}

% ─────────────────────────────────────────────────────────────────────────────
\section{Technical Specifications}

\subsection{Pixel Art Rules}
\begin{itemize}[leftmargin=1.5em, itemsep=2pt]
  \item \textbf{No anti-aliasing} --- hard pixel edges only
  \item Outline: 1px dark outline using shadow variant of each base color.
        \textbf{Never pure black \#000000}
  \item Light source: always upper left for all assets
  \item Max 20 colors per character sprite (all animation frames combined)
  \item Max 32 colors per background layer
\end{itemize}

\subsection{Owlorumo --- Display Rules}

\begin{tabularx}{\textwidth}{@{}l c c >{\raggedright\arraybackslash}X@{}}
  \toprule
  \rowcolor{tabhead}
  \color{white}\sffamily\bfseries Context &
  \color{white}\sffamily\bfseries Master &
  \color{white}\sffamily\bfseries Display &
  \color{white}\sffamily\bfseries Scale \\
  \midrule
  \rowcolor{tabrow}PassageView (gameplay) & 64x64 & 192x192 & 3x integer \\
  Library / menu & 64x64 & 256x256 & 4x integer \\
  \rowcolor{tabrow}Celebration / reward & 64x64 & 256--320+ & 4--5x integer \\
  \bottomrule
\end{tabularx}

\begin{notebox}
  Both 3x and 4x are integer scales from 64px. Godot renders pixel-perfect with
  \texttt{texture\_filter = NEAREST}. \textbf{One atlas per state, multiple display contexts.}
\end{notebox}

\subsection{Atlas Delivery Format}
\begin{itemize}[leftmargin=1.5em, itemsep=2pt]
  \item 5 PNG files: \texttt{atlas\_owlorumo\_state0.png} through \texttt{state4.png}
  \item Identical frame layout across all 5 atlases --- transparent background
  \item Frames left-to-right in rows, consistent grid, no padding
  \item Atlas max size: 2048x2048px
  \item Anchor point: center-bottom (character feet)
  \item Deliver Aseprite \texttt{.ase} source + flattened PNG
\end{itemize}

% ─────────────────────────────────────────────────────────────────────────────
\section{Scope \& Phases}

\subsection{Phase 1 --- M1 Final Slice}

Minimum assets for the first approved artist pass on top of the stable technical slice.

\begin{itemize}[leftmargin=1.5em, itemsep=2pt]
  \item[\textbf{!}] \texttt{atlas\_owlorumo\_state0.png} --- all 4 animations --- cracked crystal (\hex{2C1A0E})
  \item[\textbf{!}] \texttt{atlas\_owlorumo\_state1.png} --- all 4 animations --- faint crystal glow
  \item[\textbf{!}] \texttt{bg\_library\_layer\_4} + \texttt{layer\_3} --- far and mid shelves
  \item[\textbf{!}] \texttt{parchment\_tile} --- PassageView reading panel
  \item[\textbf{!}] \texttt{fragment\_light\_01} + \texttt{fragment\_light\_02}
  \item[\textbf{!}] \texttt{slot\_word\_empty} + \texttt{slot\_word\_filled}
\end{itemize}

\begin{notebox}
  Deliver \textbf{state\_0} before state\_1. state\_0 is the most important visual state
  --- it establishes the narrative of diminishment that the child's reading will restore.
  This is Owlorumo as the child first meets him.
\end{notebox}

\subsection{Phase 2 --- Expanded Production}
\begin{itemize}[leftmargin=1.5em, itemsep=2pt]
  \item \texttt{atlas\_owlorumo\_state2.png} + \texttt{state3.png} --- crystal lit, runes appear
  \item \texttt{bg\_library\_layer\_2} + \texttt{layer\_1} --- foreground + parallax
  \item \texttt{wood\_tile}, \texttt{tile\_word\_*}, \texttt{icon\_hint}, \texttt{btn\_back}
  \item \texttt{book\_available} + \texttt{book\_dormant}
\end{itemize}

\subsection{Phase 3 --- Late Content / M2+}
\begin{itemize}[leftmargin=1.5em, itemsep=2pt]
  \item \texttt{atlas\_owlorumo\_state4.png} --- blue mantle illuminated white-gold, \hex{D4A017} runes, identity anchors preserved (fulfilled, not replaced)
  \item Remaining animations for all states (excited, celebrate, reading, sleepy, wave)
  \item Celebration, Journal, Settings backgrounds
  \item Scramble/WordHunt tile backgrounds (16×16px NinePatchRect, no letters).
    Letters are rendered dynamically by Godot using Nunito — no letter sprites needed.
\end{itemize}

% ─────────────────────────────────────────────────────────────────────────────
\section{Working Agreement}

\subsection{Reference Images}
\begin{itemize}[leftmargin=1.5em, itemsep=2pt]
  \item Owlorumo concept art --- official character reference (embedded)
  \item Owlorumo pixel art --- AI-generated style reference (embedded)
  \item Library Scene --- AI-generated scene reference (embedded)
  \item Bookshelf close-up --- AI-generated detail reference (embedded)
  \item Parchment frame --- AI-generated PassageView reference (embedded)
\end{itemize}

\begin{notebox}
  AI references are style guides, not final designs.
  \textbf{The concept art of Owlorumo is the only non-negotiable reference.}
\end{notebox}

\subsection{Process}
\begin{itemize}[leftmargin=1.5em, itemsep=2pt]
  \item Deliver Phase 1 as rough sketches first for approval before final pixel art
  \item One revision round included per asset
  \item Deliver Aseprite \texttt{.ase} + flattened PNG
  \item All source files licensed to The Readcraftery Project upon delivery
\end{itemize}

\subsection{License}

All delivered assets become the exclusive property of The Readcraftery Project upon
full payment. The artist retains the right to show the work in their portfolio with credit.

\vspace{1cm}
\hrule
\vspace{6pt}
\begin{center}
  {\small\sffamily\color{neutralcol}
    READCRAFTERY --- Artist Bible v1.2 --- March 2026 ---
    \href{https://github.com/The-Readcraftery-Project}{github.com/The-Readcraftery-Project}
  }
\end{center}

\end{document}
